% Auto-generated from book/00_preface.md — do not edit by hand
% Regenerated by scripts/md_to_latex.py

\chapter*{Preface}
\label{ch:preface}
\addcontentsline{toc}{chapter}{Preface}
\markboth{Preface}{}

\textbf{Kingdom Come: A Quaternionic Geometric Approach to Number Theory and Physics} \emph{(Short name: QGA)}

\textbf{Mission.} Lift Hatcher's visual, diagrammatic number theory from the Farey plane and binary quadratic forms to unit quaternions, the Hopf fibration, and gauged flux lattices---keeping pure geometry and arithmetic defensible, and isolating physical models and observational hypotheses so they never masquerade as theorems.

\bigskip
\noindent\rule{\textwidth}{0.4pt}
\bigskip

\section{Why this book exists}
\label{ch:preface:why-this-book-exists}

Allen Hatcher's \emph{Topology of Numbers} (\href{https://pi.math.cornell.edu/~hatcher/TN/TNbook.pdf}{free PDF}; AMS paperback, 2022) teaches elementary number theory by \textbf{drawing}. Mediants, Farey edges, continued-fraction zigzags, and Conway topographs of binary quadratic forms turn dry Diophantine statements into spatial relationships you can almost walk along. Symmetries appear as linear fractional transformations; classification appears as finite lists of reduced forms; arithmetic depth appears as class groups of quadratic forms and fields. The method is as important as the theorems: \textbf{make the integers visible}.

\href{https://github.com/kinaar8340/kingdom_come}{Kingdom Come} is a Hopf-fibration portal---code, figures, and a Gradio laboratory---built around a different but kindred intuition. Unit quaternions form the three-sphere \(S^3\). The Hopf map \(S^3 \to S^2\) fibers that sphere into linked circles. A \textbf{gauged Hopf lattice} discretizes the geometry; \textbf{flux flywheels} are topologically protected rotating configurations on that lattice. A map from atomic number \(Z\) to flywheel metrics produces a periodic-table proxy (`\texttt{Magic Islands''). Across several observational domains a numerical signature \textbackslash{}(350/\textbackslash{}pi \textbackslash{}approx 111.41\textbackslash{}) appears as a working topological clock \textbackslash{}(W\_g\textbackslash{}). The same }`draw it, rotate it, see the links'' ethos that animates Hatcher's diagrams drives the portal's visualizers.

This book, \textbf{Kingdom Come: A Quaternionic Geometric Approach to Number Theory and Physics}, is the \textbf{lift} of Hatcher's method:

\begin{center}
\small
\begin{tabular}{@{}lp{0.55\textwidth}@{}}
\toprule
Hatcher's plane & This book's stage \\
\midrule
Farey diagram on the hyperbolic / rational plane & Discrete structure on \(S^3\), projected by Hopf \\
Binary quadratic forms & Quaternion norms, quaternary forms, \textbf{flux topographs} \\
\(SL(2,\mathbb{Z})\)-type symmetries & Left/right unit-quaternion multiplications; lattice gauge actions \\
Representations by forms & Representations by norms \textbf{and} the \(Z\mapsto\) flywheel map \\
Class groups of quadratic forms/fields & Ideal classes of quaternion orders and algebras \\
\bottomrule
\end{tabular}
\end{center}


The pure mathematics spine (Parts I--IV) aims to be defensible geometric number theory and topology. Part V states \textbf{models} and \textbf{hypotheses} about emergent physics and observational constants; those claims are labeled so they never masquerade as theorems.

\section{What this book is not}
\label{ch:preface:what-this-book-is-not}

It is \textbf{not} a modified reprint of Hatcher's PDF, nor a substitute for it. Hatcher remains the canonical geometric introduction to Farey geometry, topographs, and class groups of binary quadratic forms. We cite him freely, invite parallel reading, and---when a construction has a clear TN counterpart---say so explicitly (see \texttt{HATCHER\_MAP.md} in the project root).

It is also \textbf{not} a claim that every numerical coincidence listed in Kingdom Come is a law of nature. Where the text says \textbf{Hypothesis}, treat it as a research prompt with an eventual validation or falsification protocol (Chapter 10).

\section{Companion software and figures}
\label{ch:preface:companion-software-and-figures}

Interactive experiments live in the Kingdom Come repository and Hugging Face Space. Pedagogical lattice/topograph/composition/algebra/validation helpers live in this book's \texttt{lib/} package. Static figures are under \texttt{book/figures/} (see attribution there).

\begin{center}
\small
\begin{tabular}{@{}lp{0.55\textwidth}@{}}
\toprule
Resource & Location \\
\midrule
Book manuscript (this repo) & \texttt{\textasciitilde{}/Projects/qga/} \(\cdot\) \href{https://github.com/kinaar8340/qga}{github.com/kinaar8340/qga} \\
Kingdom Come source / portal & \texttt{\textasciitilde{}/Projects/kingdom/} \(\cdot\) \href{https://github.com/kinaar8340/kingdom_come}{github.com/kinaar8340/kingdom\_come} \\
Shared Hopf / quaternion core & \texttt{flux\_hopf\_lib} \\
Figures & \texttt{book/figures/} \\
How to run labs & [\texttt{book/HOW\_TO\_USE.md}](HOW\_TO\_USE.md) \\
\bottomrule
\end{tabular}
\end{center}


Module pointers at the end of each chapter name the Python entry points used in labs (\texttt{kingdom.core.*}, \texttt{lib.*}, and figure generators under \texttt{scripts/}).

\section{How to read}
\label{ch:preface:how-to-read}

\begin{itemize}
\item \textbf{Pure geometry and arithmetic.} Chapters 0--9; treat Chapter 10 as optional (observations and validation).  
\item \textbf{Model / portal narrative first.} Chapter 0 \(\rightarrow\) Chapter 2 \(\rightarrow\) Chapter 3 (lattice and flywheels) \(\rightarrow\) Chapter 7 \(\rightarrow\) Chapter 10.  
\item \textbf{Coming from Hatcher.} Keep TN open beside this book; use root \texttt{HATCHER\_MAP.md} as a parallel syllabus.  
\item \textbf{Coming from the Gradio app.} The Preview chapter is the `\texttt{Home + Hopf Visualizer + Flux Flywheel'' tour in prose; see also }HOW\_TO\_USE.md`.  
\item \textbf{Open problems.} Living list in \texttt{notes/open\_problems.md} (OP1--OP6).
\end{itemize}

Claim labels used throughout (see also \texttt{HOW\_TO\_USE.md}):

\begin{center}
\small
\begin{tabular}{@{}lp{0.55\textwidth}@{}}
\toprule
Label & Meaning \\
\midrule
\textbf{Theorem} & Proved here or classical \\
\textbf{Model} & Consistent mathematical or physical construction, not claimed as nature's unique law \\
\textbf{Hypothesis} & Observational claim needing validation or falsification (Chapter 10 / Table T4) \\
\textbf{Software fact} & True of the current code or portal implementation \\
\bottomrule
\end{tabular}
\end{center}


\section{Acknowledgments}
\label{ch:preface:acknowledgments}

Allen Hatcher's freely available \emph{Topology of Numbers} made geometric elementary number theory a shared public good; this project stands on that invitation to \emph{see} number theory. The Kingdom Come codebase and observation program (Hopf visualizers, lattice simulations, flux flywheel element maps, Magic Island sweeps, and multi-domain \(350/\pi\) investigations) supply the computational and visual backbone of the lift.

\bigskip
\noindent\rule{\textwidth}{0.4pt}
\bigskip


